\documentclass[11pt]{article}

\usepackage[a4paper,margin=1in]{geometry}
\usepackage{booktabs}
\usepackage{listings}
\usepackage{siunitx}
\usepackage{hyperref}

\title{Problem 2: OpenMP Acceleration and Correctness Test for the NEP Force Path}
\author{}
\date{May 23, 2026}

\begin{document}
\maketitle

\section{Goal}

Problem 2 requires OpenMP acceleration for molecular-dynamics potential
calculation, with special attention to data dependencies, reductions, and
correctness under different thread counts.  In this work, the target is the NEP
CPU implementation under \texttt{source/source\_md/potential/ml/nep}.

\section{Implementation}

The descriptor and neighbor-list construction paths already contained OpenMP
parallel loops.  The remaining risky hotspot was the force path:
\begin{itemize}
  \item \texttt{find\_force\_radial\_small\_box}
  \item \texttt{find\_force\_angular\_small\_box}
  \item \texttt{find\_force\_ZBL\_small\_box}
\end{itemize}

These loops cannot be parallelized safely with a plain
\texttt{\#pragma omp parallel for}, because each neighbor contribution updates
both atoms in a pair:
\[
  F_{n1} \mathrel{+}= f_{12}, \qquad F_{n2} \mathrel{-}= f_{12}.
\]
Different OpenMP threads can therefore write to the same force and virial
entries at the same time.

The implemented solution uses thread-local accumulation buffers:
\begin{enumerate}
  \item Each OpenMP thread allocates private force and virial arrays.
  \item The outer atom loop is parallelized with \texttt{\#pragma omp for}.
  \item Pair contributions are accumulated only into the calling thread's
        private arrays.
  \item After the parallel loop, private arrays are reduced into the shared
        output arrays inside a short critical section.
\end{enumerate}

This keeps the force calculation race-free while preserving the original
pairwise force and virial formulas.

\section{Correctness Test}

A small standalone NEP checker was added:
\begin{itemize}
  \item Source: \texttt{source/source\_md/tools/nep\_omp\_check.cpp}
  \item Runner: \texttt{source/source\_md/tools/run\_nep\_omp\_check.sh}
  \item Model: \texttt{tests/PP\_ORB/nep\_hfo2.txt}
\end{itemize}

The checker computes the same six-atom Hf/O configuration using one OpenMP
thread as the reference, then repeats the calculation with multiple thread
counts.  It reports average runtime and maximum absolute differences for
potential, force, and virial arrays.

\section{Initial Run Results}

The following command was used on the local machine:

\begin{lstlisting}[language=bash]
tools/run_nep_omp_check.sh ../../tests/PP_ORB/nep_hfo2.txt 1,2,4 100
\end{lstlisting}

The compiler selected by the script was \texttt{g++-15} with
\texttt{-O2 -fopenmp}.  The results were:

\begin{table}[h]
\centering
\begin{tabular}{@{}rrrrrr@{}}
\toprule
Threads & Avg. time (ms) & Speedup & Max $\Delta E$ & Max $\Delta F$ & Max $\Delta V$ \\
\midrule
1 & 0.96246 & 1.00 & 0 & 0 & 0 \\
2 & 0.78283 & 1.23 & 0 & $2.22045\times10^{-16}$ & $8.88178\times10^{-16}$ \\
4 & 0.80320 & 1.20 & 0 & $4.44089\times10^{-16}$ & $1.33227\times10^{-15}$ \\
\bottomrule
\end{tabular}
\caption{Initial NEP OpenMP correctness and runtime check.}
\end{table}

\section{Analysis}

The force and virial differences are at floating-point roundoff level, so the
OpenMP force path is numerically consistent with the single-thread reference.
The two-thread and four-thread runs show modest speedups of about $1.23\times$
and $1.20\times$ for this small test case.  The four-thread result is slightly
slower than the two-thread result, which is reasonable for such a small system
because thread startup, private buffer allocation, and final reduction overhead
can dominate the actual NEP work.

For larger systems, the per-atom NEP descriptor and force work should become
large enough for the OpenMP parallel region to amortize its overhead more
effectively.  Therefore, the current result should be interpreted mainly as a
correctness and initial-performance validation, not as the final scalability
limit.

\section{How to Reproduce}

\begin{lstlisting}[language=bash]
cd source/source_md
tools/run_nep_omp_check.sh
\end{lstlisting}

Optional arguments are:

\begin{lstlisting}[language=bash]
tools/run_nep_omp_check.sh MODEL_FILE THREAD_LIST REPEATS
tools/run_nep_omp_check.sh /path/to/nep.txt 1,2,4,8 50
\end{lstlisting}

\end{document}
